\section{Hebrew Sequence and Vision}

The identified 1909 Kittel Hebrew printing uses wayyiqtol forms to advance the
represented actions at verse 5. They supply narrated progression, not a measure
of elapsed time, and cannot by themselves rule out a gap. At verse 12, the Qal
infinitive construct with prefixed \emph{lamed} presents the sun approaching
its setting. At verse 17, the feminine singular Qal active participle agrees
with ``the sun''; it is not a completed verbal form. Together with the thick
darkness and verse 18's ``in that day,'' the clauses create a clock inside the
sacrifice scene. The question is whether that clock also governs verse 5.

The vision notice in verse 1 prevents narrated sequence from being identified
without argument with ordinary waking chronology. It does not erase sequence
inside the represented scene; it changes what kind of sequence may be in view.
The bound Augustine witness calls the episode a vision without using that fact
to identify the hour of verse 5.

